% !TeX root = ../../main.tex
\subsection{反比例}

\begin{definitionbox}{\textbf{反比例}}
    \label{def:反比例}
    $y$ が $x$ の関数で，$x$ と $y$ の関係が次のような式で表されるとき、\textbf{$y$ は $x$ に反比例する}という.

    \[\fitblankbf[]{y = \frac{a}{x}} \qquad (x \neq 0)\]

    \vspace{0.5em}
    このとき、$a$ を\fitblankbf[]{比例定数} という.${}^{*}$
\end{definitionbox}

反比例は、一方の数量が2倍、3倍 … となると、それに伴ってもう一方の数量は、$\dfrac{1}{2}$ 倍、$\dfrac{1}{3}$ 倍 … となるような関係をいう。

\begin{theorembox}{\textbf{反比例のグラフ}}
    反比例：$y = \dfrac{a}{x}$ のグラフは \fitblankbf[]{双曲線}である.
    \begin{enumerate}
        \item $a > 0$ のとき,\fitblankbf[]{右上} と \fitblankbf[]{左下} にある
        \item $a < 0$ のとき,\fitblankbf[]{右下} と \fitblankbf[]{左上} にある
        \item $x$軸 と $y$軸 には \fitblankbf[]{交わらない}
    \end{enumerate}
\end{theorembox}

\begin{example}[3]
    $y$ が $x$ に反比例し，$x = 4$ のとき $y = 3$ である. $y$ を $x$ の式で表しなさい.
\end{example}

\begin{proof}\mbox{}\\
    $y = \dfrac{a}{x}$ とおく.

    \begin{hiddenanswer}
        $x = 4$, $y = 3$ を代入すると，
        \begin{align*}
            3 &= \frac{a}{4} \\
            a &= 12
        \end{align*}
        \finalanswer{答え：$y = \dfrac{12}{x}$}
    \end{hiddenanswer}
\end{proof}

\begin{example}[4]
    $y$ が $x$ に反比例し，$x = 2$ のとき $y = 6$ である. $x = -3$ のときの $y$ の値を求めなさい.
\end{example}

\begin{proof}\mbox{}\\
    $y = \dfrac{a}{x}$ とおく.

    \begin{hiddenanswer}
        $x = 2$, $y = 6$ を代入すると，
        \begin{align*}
            6 &= \frac{a}{2} \qquad \Longleftrightarrow \qquad a = 12
        \end{align*}

        $x = -3$ のときの $y$ の値を求めると，
        \begin{align*}
            y &= \frac{12}{-3} \\
            y &= -4
        \end{align*}
        \finalanswer{答え：$y = -4$}
    \end{hiddenanswer}
\end{proof}

\vfill
\noindent\rule{\textwidth}{0.4pt}

\vspace{0.2em}
\noindent
$^{*}$ Definition \ref{def:反比例} より、$x \neq 0$ のとき、$\bm{a = xy}$ となる.

\newpage
